\section{Related Works} \label{sec:02-related}
\subsection{Text Dataset Distillation} \label{sec:02-related-textdd}
\begin{foldable} % DD --> Text DD
  Dataset distillation was introduced by~\cite{wang2018dataset} as a bi-level meta-learning problem: synthesize a small set of images such that a model trained on them matches one trained on the full dataset.
  Subsequent methods improve scalability and fidelity via gradient matching~\cite{zhao2020dataset}, distribution matching~\cite{wang2022cafe}, and trajectory matching~\cite{cazenavette2022dataset,kim2022dataset}.
  These methods share a common assumption: data is continuous and differentiable, so the synthetic set can be directly optimised via backpropagation.
  Text violates the differentiability assumption: token decoding is non-differentiable, outputs must be semantically coherent, and embedding spaces are tightly coupled to specific architectures.
\end{foldable}

\begin{foldable} % Text DD landscape: progression of trade-offs and the uniform-weighting gap
  Adapting DD to text has therefore required fundamentally different strategies, from meta-learning on synthetic embeddings to LLM-based generation of readable corpora, each generation trading one limitation for another.
  SLDD and DDTC~\cite{sucholutsky2021soft,li2021data} directly adapted the meta-learning formulation~\cite{wang2018dataset}, optimizing synthetic text embeddings so that a model trained on them minimizes task loss.
  Of the two, only SLDD partially recovers interpretability by mapping distilled embeddings to per-token nearest bag-of-words.
  DDAL~\cite{maekawa2023dataset} optimizes the attention labels of a BERT model to match class-wise dataset statistics, achieving partial task alignment but producing outputs that are both architecture-locked and uninterpretable.
  A second generation of methods addressed interpretability by rethinking the pipeline entirely.
  DiLM~\cite{maekawa2024dilm} generates readable candidate texts via an LLM and selects among them by gradient matching over the training corpus, followed by k-centre selection; but matching gradients at discrete checkpoints is a proxy for task loss, leaving task alignment only partial.
  DaLLME~\cite{tao2024textual} learns an inverse mapping from distilled embeddings back to synthetic text, recovering interpretability and attempting global coverage via k-centroids clustering; but with no task-loss objective, task alignment is absent.
  Table~\ref{tab:landscape-textdd-methods} maps this progression: each generation addressed some gaps while leaving others open.
\end{foldable}

\begin{table}[ht] % Table: Comparison of text DD methods.
  \centering
  \caption{Landscape of text DD methods.}
  \label{tab:landscape-textdd-methods}
  \setlength{\tabcolsep}{8pt}
  \begin{tabular}{l c c c c c c}
    \toprule
    \textbf{Gap}         & SLDD         & DDTC         & DDAL         & DiLM         & DaLLME       & \textbf{TAKE} \\
    \midrule
    Task alignment       & $\checkmark$ & $\checkmark$ & $\triangle$  & $\triangle$  & $-$          & $\checkmark$ \\
    Global optimization  & $\triangle$  & $\triangle$  & $\triangle$  & $\triangle$  & $\triangle$  & $\checkmark$ \\
    Interpretability     & $\triangle$  & $-$          & $-$          & $\checkmark$ & $\checkmark$ & $\checkmark$ \\
    Importance weighting & $-$          & $-$          & $-$          & $-$          & $-$          & $\checkmark$ \\
    \bottomrule
    \multicolumn{7}{l}{$-$: absent, $\triangle$: partial, $\checkmark$: addressed}
  \end{tabular}
\end{table}

\begin{foldable} % Uniform-weighting gap: the blind spot shared by all methods
  A subtler gap cuts across all methods: every existing approach applies \emph{uniform weighting}, treating every training sample as equally informative.
  At moderate distillation budgets this is defensible, but extreme compression is precisely the regime where distillation matters most, and there the assumption fails.
  Samples vary substantially in their influence on learning (clean vs.\ noisy, easy vs.\ hard), yet no prior method accounts for this structure.
  Without importance weighting, the distillation budget is exhausted on uninformative samples, biasing selection away from the boundary cases that matter most for robustness.
\end{foldable}

\subsection{Influence Functions and the Hard-Sample Bias} \label{sec:02-related-influence}
\begin{foldable} % Influence functions & trajectory-based methods
  Influence functions (IF)~\cite{cook1982residuals,koh2017understanding} quantify the leave-one-out effect of a training point via $\mathcal{I}(z) = \nabla^\top H_\theta^{-1} \nabla$, but depend on a single checkpoint and require expensive second-order computation.
  Trajectory-based methods~\cite{pruthi2020estimating,park2023trak,kwon2023datainf} extend attribution across checkpoints but are \emph{test-conditioned}: designed to attribute a fixed test prediction.
  The self-influence approximation adapts them to corpus-level scoring by using each training point as its own query.
  However, self-influence is evaluated only once at the final parameters, introducing a systematic bias that corrupts the corpus-level scores.
\end{foldable}

\begin{foldable} % Hard-sample bias
  We refer to this as the \emph{hard-sample bias} (HSB): at convergence, clean samples have near-zero gradient by definition, so noisy or hard samples dominate influence scores.
  This failure mode is well-documented in adjacent work: \cite{arpit2017closer} show that networks fit clean patterns before noisy ones; dataset cartography~\cite{swayamdipta2020dataset} confirms that single-checkpoint statistics conflate difficulty with noise; and Co-teaching~\cite{han2018co} documents systematic overselection of noisy samples in gradient-based curricula.
  In standard training, HSB can often be mitigated via regularization; in the extreme low-data regime of DD, however, it amplifies: biased selection suppresses the \emph{informative} samples essential for model robustness.
  To our knowledge, no prior text DD method has identified or corrected for it.
\end{foldable}

\subsection{Optimal Transport for Distribution Matching} \label{sec:02-related-ot}
\begin{foldable} % Distribution matching as a DD surrogate
  Distribution matching has been proposed as a tractable DD surrogate~\cite{zhao2023dataset}: if the distilled and training sets induce equal expected gradients along the optimization trajectory, the trained models converge to comparable parameters.
  This is a sufficient condition, but not an equivalence (the formal gap bound appears in Theorem~\ref{thm:gap}).
  Common divergences have practical shortcomings at the extreme budgets relevant here.
  MMD~\cite{gretton2012kernel} lacks sample-level assignment structure, requiring auxiliary scoring heuristics that reintroduce the global optimization gap that distribution matching is meant to close.
  Kullback--Leibler and Jensen--Shannon divergences require density estimation, which is unreliable in the low-data regime~\cite{paninski2003estimation}.
\end{foldable}

\begin{foldable} % Discrete OT
  OT in the Monge-Kantorovich sense avoids both pitfalls: it operates natively on discrete measures, respects the metric geometry of the embedding space, and produces an explicit transport plan that directly assigns training samples to prototypes---no auxiliary heuristic needed.
  Entropic regularization via Sinkhorn-Knopp~\cite{cuturi2013sinkhorn} makes this tractable at corpus scale.
  OT has been widely applied in NLP and machine learning~\cite{kusner2015word,tolstikhin2018wasserstein,courty2016optimal,arjovsky2017wasserstein}, but never to text DD, and never combined with a knowledge-reweighted source distribution in any DD setting.
  TAKE fills both gaps: $\kappa_n$ defines a non-uniform source measure, and the discrete entropic OT plan selects prototypes covering the knowledge-weighted training distribution without any auxiliary selection step.
\end{foldable}
